\chapter{Historiographie et état de la recherche}\label{ch:historiographie}

\section{Méthode et limites de la recherche bibliographique}

L’état de la recherche repose sur des catalogues de bibliothèques et des ressources éditoriales : catalogues universitaires (Stanford, Berkeley, La Sapienza, RERO), HathiTrust, Internet Archive, Persée, Crossref, PubMed, Cambridge Core, sites des éditeurs (Oxford University Press, De Gruyter, École française de Rome, Éditions de l’EHESS, Pontificio Istituto di Archeologia Cristiana) et comptes rendus de la \emph{Bryn Mawr Classical Review}. Cette recherche établit l’existence et les données bibliographiques des travaux cités. Elle ne remplace pas leur lecture intégrale. Lorsqu’une position est attribuée à un auteur d’après un compte rendu ou une présentation d’éditeur, la note le dit. Lorsque les sources divergent sur une pagination ou une date, la divergence est signalée. La bibliographie finale (chapitre~\ref{ch:bibliographie}) indique le niveau de vérification de chaque titre.

La couverture privilégie les instruments (recueils, fastes, prosopographies), les synthèses et les travaux qui ont structuré les débats dont relève le corpus : la datation des inscriptions chrétiennes de Rome, l’histoire du formulaire, l’Église romaine et l’aristocratie, la sépulture et le culte des saints, l’usage démographique et quantitatif des inscriptions, l’épigraphie numérique. Elle est plus faible pour la littérature italienne récente consacrée à chaque cimetière et pour les études ponctuelles d’inscriptions, qui supposent l’accès aux bibliothèques spécialisées de Rome. Aucune étude spécifiquement consacrée aux formules trinitaires ou christologiques des épitaphes romaines n’a été identifiée ; ce constat ne prouve pas qu’il n’en existe pas.

\section{De Rossi : la date comme fondement}

L’épigraphie chrétienne de Rome a reçu son premier cadre méthodique de Giovanni Battista de Rossi. Le premier volume des \emph{Inscriptiones christianae urbis Romae septimo saeculo antiquiores}, imprimé de 1857 à 1861, réunit les inscriptions pourvues d’une date explicite. Le second (II, 1, 1888) est consacré aux sylloges médiévales, qui ont transmis une partie des textes perdus. Une première partie de supplément, due à G. Gatti, a paru en 1915\footnote{G. B. de Rossi, \emph{Inscriptiones christianae urbis Romae septimo saeculo antiquiores}, I, Rome, 1857–1861 ; II, 1, Rome, 1888 (vérifié : HathiTrust, Internet Archive). Supplément de G. Gatti, 1915 : notices de catalogue.}. Le choix de publier d’abord les inscriptions datées fonde une méthode : dater les inscriptions sans date par comparaison avec celles qui en ont une. L’histoire du formulaire, des abréviations, des symboles et des écritures s’est construite sur ce socle.

Cette méthode comporte un risque de circularité. Si des inscriptions datées par leur formulaire sont réintroduites dans un corpus servant à décrire l’évolution de ce formulaire, la chronologie se confirme elle-même. Le présent corpus exclut en principe les inscriptions datées par des critères paléographiques, formulaires ou topographiques. Il reste que l’année de quelques textes repose sur une restitution de l’éditeur ou sur la seule base de données (chapitre~\ref{ch:fondement}) : la frontière entre date transmise et date inférée passe à l’intérieur même du corpus. La représentativité des inscriptions datées pour dater les autres reste une question ouverte de la discipline. Le volume n’y répond pas, mais il en documente les conditions pour une période.

\section{La nouvelle série des ICUR : de la chronologie à la topographie}

La \emph{nova series} des ICUR, inaugurée par Angelo Silvagni en 1922 et dont le dixième volume a paru en 1992, a substitué à l’ordre chronologique de De Rossi un ordre topographique. Le premier volume réunit les inscriptions d’origine incertaine ; les neuf suivants sont organisés par voies et par groupes de cimetières : II, voies Cornelia, Aurelia, Portuensis et d’Ostie (Silvagni, 1935) ; III, voie Ardéatine (Silvagni et Ferrua, 1956) ; IV, entre les voies Appienne et Ardéatine (Ferrua, 1964) ; V, voie Appienne (Ferrua, 1971) ; VI, voies Latine, Labicane et Prénestine (Ferrua, 1975) ; VII, voie Tiburtine (Ferrua, 1980) ; VIII, voies Nomentane et Salaria (Ferrua, 1983) ; IX, autres cimetières de la Salaria (Ferrua et Mazzoleni, 1985) ; X, Salaria vetus et voie Flaminienne (Mazzoleni et Carletti, 1992)\footnote{Composition vérifiée sur les catalogues de La Sapienza, de Stanford et de RERO et sur la liste des publications du Pontificio Istituto di Archeologia Cristiana. Certaines bibliographies inversent l’ordre des éditeurs du volume X ; les catalogues donnent « Mazzoleni, Carletti ».}.

Ce changement de principe est décisif pour l’historien. Il rattache chaque inscription à un secteur funéraire quand celui-ci est connu ; mais il disperse les inscriptions datées dans dix volumes et concentre dans le premier celles dont la provenance s’est perdue. Le corpus reflète cette structure : le volume I fournit la première référence de \st{nVolI} des \st{nC} textes funéraires. L’Epigraphic Database Bari, qui suit l’ordre des ICUR et en précise souvent les rattachements, permet aujourd’hui d’aller au-delà du volume et de nommer le cimetière pour une large part du corpus (chapitre~\ref{ch:provenance}).

\section{Diehl et les instruments de concordance}

Le recueil d’Ernst Diehl, \emph{Inscriptiones Latinae Christianae Veteres} (ILCV), publié à Berlin de 1925 à 1931, réimprimé en 1961 et complété en 1967 par un \emph{Supplementum} de J. Moreau et H.-I. Marrou publié à Dublin et Zurich, organise les inscriptions chrétiennes latines selon une logique thématique — fonctions, formules, catégories de défunts — et non topographique\footnote{E. Diehl, \emph{Inscriptiones Latinae Christianae Veteres}, I–III, Berlin, Weidmann, 1925–1931 ; réimpr. 1961 ; IV, \emph{Supplementum}, éd. J. Moreau, H.-I. Marrou, Dublin – Zurich, Weidmann, 1967 (vérifié : HathiTrust, Stanford SearchWorks).}. C’est l’instrument qui a rendu possible l’étude comparée du formulaire. Dans les concordances transmises par l’EDCS, certaines références à l’ILCV sont précédées du signe « + », et plusieurs paires de notices du corpus partagent une telle référence sans être des doublons. La valeur de ce signe n’a pas été vérifiée dans la documentation de la base ni dans la préface du recueil : il n’est utilisé ici que comme indice de rapprochement, non comme preuve d’identité.

\section{Les synthèses sur le formulaire et l’inscription funéraire}

L’étude du formulaire funéraire chrétien s’est développée dans une série de manuels et de synthèses : le traité de F. Grossi Gondi (1920), l’\emph{Archeologia cristiana} de P. Testini (1958 ; 2\textsuperscript{e}~éd. 1980), le manuel de D. Mazzoleni (2002), l’anthologie commentée de C. Carletti, \emph{Iscrizioni cristiane di Roma} (1986), et sa synthèse \emph{Epigrafia dei cristiani in Occidente dal III al VII secolo. Ideologia e prassi} (Bari, 2008)\footnote{Notices vérifiées : Grossi Gondi 1920 (catalogues en ligne) ; Testini (Zenon 000158599) ; Mazzoleni 2002 (Zenon 000150908 ; Rome, Pontificia Università Lateranense) ; Carletti 1986 (guide heuristique de l’université de Gand : anthologie de 142 inscriptions, surtout tirées des ICUR) ; Carletti 2008 (bibliographie de la base \emph{The Cult of Saints in Late Antiquity}, Oxford).}. Le sous-titre de ce dernier ouvrage résume l’orientation de la recherche récente : le formulaire y est étudié comme le produit d’une idéologie et d’une pratique, c’est-à-dire d’un système de choix contraint par la tradition, la technique et le coût, autant que comme l’expression d’une croyance. Le présent volume retient cette orientation sans prétendre en restituer les analyses.

Dans le domaine francophone, la synthèse de référence est l’article « Grabinschrift II (lateinisch) » de Charles Pietri, publié dans le \emph{Reallexikon für Antike und Christentum} (t. 12, 1983, col. 514–590), dont la version française originale, « Inscriptions funéraires latines », a été reprise dans \emph{Christiana respublica} (1997, p. 1407–1468)\footnote{Ch. Pietri, « Grabinschrift II (lateinisch) », \emph{Reallexikon für Antike und Christentum}, 12, Stuttgart, 1983, col. 514–590 ; version française dans \emph{Christiana respublica. Éléments d’une enquête sur le christianisme antique}, Rome, 1997 (Collection de l’École française de Rome, 234), p. 1407–1468 (vérifié ; une bibliographie de Cambridge donne fautivement « 415–590 »).}. Selon les éditeurs du volume X des ICUR, cités par la base \emph{The Cult of Saints}, la date de déposition, absente des épitaphes de martyrs du \textsc{iii}\textsuperscript{e}~siècle, n’apparaît qu’au début du \textsc{iv}\textsuperscript{e}\footnote{Base \emph{The Cult of Saints in Late Antiquity} (université d’Oxford), notice E07620 ; information de seconde main.}. Le corpus étudié appartient donc à une époque où la date de déposition, liée à la commémoration de l’anniversaire, est déjà un élément du formulaire.

\section{Charles Pietri, l’Église de Rome et la société chrétienne}

L’œuvre de Charles Pietri domine l’histoire institutionnelle et sociale de la Rome chrétienne du \textsc{iv}\textsuperscript{e}~siècle. Sa thèse, \emph{Roma christiana. Recherches sur l’Église de Rome, son organisation, sa politique, son idéologie de Miltiade à Sixte III (311–440)}, soutenue en 1973 et publiée en 1976 en deux volumes dans la Bibliothèque des Écoles françaises d’Athènes et de Rome (n\textsuperscript{o}~224), reconstitue l’organisation de l’Église romaine, ses ressources, son clergé et l’idéologie de ses évêques\footnote{Ch. Pietri, \emph{Roma christiana…}, Rome, École française de Rome, 1976 (BEFAR, 224), 2 vol., DOI 10.3406/befar.1976.1278 (vérifié : Persée, WorldCat ; comptes rendus de H. Chadwick, \emph{Journal of Ecclesiastical History}, 30, 1979, p. 477–479 ; J.-M. Sansterre, \emph{L’Antiquité classique}, 47, 1978 ; P. Nautin, \emph{Revue de l’histoire des religions}, 198, 1981). La pagination totale varie selon les sources (1750, 1792 ou 1830 p.) ; elle n’est pas donnée ici.}. Elle est complétée par un « Appendice prosopographique » consacré au clergé romain\footnote{Ch. Pietri, « Appendice prosopographique à la \emph{Roma christiana} (311–440) », \emph{Mélanges de l’École française de Rome. Antiquité}, 89/1, 1977, p. 371–415 (vérifié : Persée).}, par les articles réunis dans \emph{Christiana respublica}, dont une étude de l’onomastique chrétienne de Rome\footnote{Ch. Pietri, « Remarques sur l’onomastique chrétienne de Rome », dans \emph{L’onomastique latine}, Paris, CNRS, 1977, p. 437–456 selon Persée (une autre source donne p. 437–445) ; repris dans \emph{Christiana respublica}, p. 1571–1582.}, et par la \emph{Prosopographie chrétienne du Bas-Empire} consacrée à l’Italie, achevée après sa mort sous la direction de Luce Pietri\footnote{Ch. Pietri (†), L. Pietri (dir.), \emph{Prosopographie chrétienne du Bas-Empire}, 2, \emph{Prosopographie de l’Italie chrétienne (313–604)}, Rome, École française de Rome, 1999–2000, 2 vol. (vérifié : Persée ; notices de libraires).}.

Pour le présent corpus, cette œuvre fournit le cadre dans lequel s’inscrivent les mentions de clercs (lecteurs de titres, prêtres, diacres), de l’évêque Sirice et des fossoyeurs. Elle fournit aussi l’instrument prosopographique qui permettra de contrôler les identifications proposées dans les dossiers (chapitre~\ref{ch:rangs}) ; ces identifications restent provisoires tant que ce contrôle n’a pas été fait.

L’étude de la christianisation de l’aristocratie occidentale a été renouvelée par M. R. Salzman, \emph{The Making of a Christian Aristocracy} (Cambridge, Mass., 2002), qui, selon le compte rendu de J. O’Donnell, repose sur l’examen de 414 aristocrates sénatoriaux d’Occident et aborde la conversion par le statut et les valeurs plutôt que par l’initiative des évêques\footnote{M. R. Salzman, \emph{The Making of a Christian Aristocracy. Social and Religious Change in the Western Roman Empire}, Cambridge (Mass.), Harvard University Press, 2002 ; J. O’Donnell, \emph{Bryn Mawr Classical Review}, 2002.06.04 (compte rendu vu en extrait).}. Le corpus ne contient qu’une dizaine de textes où une personne commémorée porte un titre de l’ordre sénatorial, une fois exclue la titulature des consuls : il documente quelques familles clarissimes, sans pouvoir mesurer la christianisation de l’aristocratie.

\section{La sépulture, le culte des saints et les fossoyeurs}

La question de la proximité des tombes des martyrs a été posée en termes généraux par Peter Brown, \emph{The Cult of the Saints} (Chicago, 1981)\footnote{P. Brown, \emph{The Cult of the Saints. Its Rise and Function in Latin Christianity}, Chicago, University of Chicago Press, 1981 (Haskell Lectures on History of Religions, n.s., 2) (partiellement vérifié : notice Internet Archive).}, et reprise pour Rome par Steffen Diefenbach, qui étudie la mémoire des saints comme lieu de construction d’identités collectives du \textsc{iii}\textsuperscript{e} au \textsc{v}\textsuperscript{e}~siècle\footnote{S. Diefenbach, \emph{Römische Erinnerungsräume. Heiligenmemoria und kollektive Identitäten im Rom des 3. bis 5. Jahrhunderts n. Chr.}, Berlin – New York, De Gruyter, 2007 (Millennium-Studien, 11), DOI 10.1515/9783110204261 (vérifié : De Gruyter ; l’ouvrage est en accès ouvert).}. La poésie épigraphique de l’évêque Damase, consacrée pour une large part aux martyrs, a fait l’objet d’une édition commentée par D. E. Trout (Oxford, 2015), qui complète celle d’A. Ferrua (1942)\footnote{A. Ferrua, \emph{Epigrammata Damasiana}, Cité du Vatican, 1942 (Sussidi allo studio delle antichità cristiane, 2) (vérifié : liste PIAC) ; D. E. Trout, \emph{Damasus of Rome. The Epigraphic Poetry}, Oxford, 2015 (Oxford Early Christian Texts) (partiellement vérifié : pagination divergente selon les notices).}.

La gestion des sépultures a été reformulée par Éric Rebillard. Selon les comptes rendus, \emph{Religion et sépulture} (2003 ; trad. angl. \emph{The Care of the Dead in Late Antiquity}, 2009) conteste la thèse, issue de De Rossi, selon laquelle l’Église aurait possédé et organisé des cimetières communautaires dès le \textsc{iii}\textsuperscript{e}~siècle, ainsi que l’hypothèse des collèges funéraires. La sépulture des fidèles y apparaît comme l’affaire des familles, l’Église ne prenant en charge que celle des martyrs, et les fossoyeurs comme des opérateurs vendant des emplacements\footnote{É. Rebillard, \emph{Religion et sépulture. L’Église, les vivants et les morts dans l’Antiquité tardive}, Paris, Éditions de l’EHESS, 2003 (Civilisations et sociétés, 115) ; \emph{The Care of the Dead in Late Antiquity}, Ithaca, Cornell University Press, 2009 (Cornell Studies in Classical Philology, 59) (vérifié). Résumé d’après les comptes rendus : \emph{Bryn Mawr Classical Review}, 2004.07.52 ; M. Costambeys, \emph{sehepunkte}, 5, 2005, n\textsuperscript{o}~11. Ce dernier rappelle que Pietri avait établi que les fossoyeurs n’étaient pas des clercs.}. Les épitaphes du corpus qui mentionnent l’achat d’un emplacement à un fossoyeur (chapitre~\ref{ch:sepultures}) entrent directement dans ce débat ; elles ne permettent pas de le trancher, mais elles fournissent des cas datés.

Les catacombes elles-mêmes, leur développement et leurs décors sont présentés dans la synthèse de V. Fiocchi Nicolai, F. Bisconti et D. Mazzoleni, publiée à Ratisbonne chez Schnell \& Steiner en 1998 et plusieurs fois rééditée\footnote{V. Fiocchi Nicolai, F. Bisconti, D. Mazzoleni, \emph{Le catacombe cristiane di Roma. Origini, sviluppo, apparati decorativi, documentazione epigrafica}, Ratisbonne, Schnell \& Steiner, 1998 (notice vérifiée ; l’éditeur est confirmé par la 4\textsuperscript{e}~édition, 2025, sous le sous-titre \emph{Storia, immagini, epigrafia}, ISBN 978-3-7954-3967-5 ; traductions et éditions intermédiaires non contrôlées).}. Les rattachements de l’EDB permettent désormais de situer plus de la moitié des textes funéraires dans cette topographie (chapitre~\ref{ch:provenance}).

\section{Les âges au décès, les saisons et la démographie}

L’abondance des âges au décès dans les épitaphes romaines a nourri des tentatives de reconstitution démographique. Keith Hopkins a montré en 1966 que ces âges ne peuvent servir à estimer l’espérance de vie : la commémoration sélectionne les défunts, et les âges déclarés sont affectés par l’arrondi\footnote{K. Hopkins, « On the Probable Age Structure of the Roman Population », \emph{Population Studies}, 20/2, 1966, p. 245–264, DOI 10.1080/00324728.1966.10406097 (vérifié : PubMed, PMID 22084913).}. Henrik Nordberg avait étudié dès 1963 les indications des inscriptions chrétiennes de Rome relatives à la durée de la vie et aux dates de naissance et de décès\footnote{H. Nordberg, \emph{Biometrical Notes. The Information on Ancient Christian Inscriptions from Rome concerning the Duration of Life and the Dates of Birth and Death}, Helsinki, 1963 (Acta Instituti Romani Finlandiae, II, 2) (notice vérifiée).}. Brent D. Shaw a déplacé l’usage de ces données vers la saisonnalité de la mortalité, à partir des dates de déposition des épitaphes chrétiennes\footnote{B. D. Shaw, « Seasons of Death: Aspects of Mortality in Imperial Rome », \emph{Journal of Roman Studies}, 86, 1996, p. 100–138, DOI 10.2307/300425 (vérifié : Cambridge Core).}. Walter Scheidel a repris et discuté la critique de Hopkins\footnote{W. Scheidel, « Roman Age Structure: Evidence and Models », \emph{Journal of Roman Studies}, 91, 2001, p. 1–26 (partiellement vérifié ; DOI vu de seconde main, non donné).}.

Le présent volume reprend ces mesures sur un corpus restreint, en distinguant les âges entièrement conservés des âges restitués et en dénombrant les personnes plutôt que les notices (chapitre~\ref{ch:personnes}). Il rapporte des indicateurs d’arrondi et une distribution mensuelle des dates (chapitre~\ref{ch:calendrier}). Il n’en tire aucune conclusion sur la mortalité.

\section{L’« habitude épigraphique » et les corpus comme données}

Ramsay MacMullen a mis en évidence les variations de la production des inscriptions dans l’Empire et la nécessité de les expliquer comme une pratique culturelle, non comme un simple reflet de la population\footnote{R. MacMullen, « The Epigraphic Habit in the Roman Empire », \emph{American Journal of Philology}, 103/3, 1982, p. 233–246, DOI 10.2307/294470 (métadonnées contrôlées sur JSTOR ; article non relu).}. Elizabeth A. Meyer a proposé de rapporter cette habitude, pour les épitaphes, à la valeur sociale du statut de citoyen\footnote{E. A. Meyer, « Explaining the Epigraphic Habit in the Roman Empire: The Evidence of Epitaphs », \emph{Journal of Roman Studies}, 80, 1990, p. 74–96, DOI 10.2307/300281 (vérifié : Cambridge Core).}. L’abondance de l’épigraphie funéraire chrétienne de Rome au \textsc{iv}\textsuperscript{e}~siècle pose à ces modèles une question que le présent volume ne traite pas, faute de données de comparaison.

La dernière décennie a vu se développer l’exploitation quantitative des grandes bases épigraphiques. P. Heřmánková, V. Kaše et A. Sobotková ont formulé, à partir des exports de l’EDH et de l’EDCS diffusés par le projet SDAM, le programme d’une épigraphie « macro-historique » fondée sur des données nettoyées et documentées\footnote{P. Heřmánková, V. Kaše, A. Sobotková, « Inscriptions as data: digital epigraphy in macro-historical perspective », \emph{Journal of Digital History}, 1/1, 2021, DOI 10.1515/jdh-2021-1004 (vérifié : Crossref ; pagination à confirmer). Dépôt de code et de données : github.com/sdam-au/EDCS\_ETL.}. Trois de leurs observations fondent ici des règles de méthode : les inscriptions conservées ne constituent pas un échantillon stratifié ; la précision apparente d’une donnée peut dépasser celle de l’information réellement documentée ; le classement religieux d’une inscription ne se substitue pas à une classification fonctionnelle. Le présent volume en fournit un cas d’étude, puisque la confrontation de deux bases y révèle des pertes que l’usage d’une seule ne laissait pas soupçonner (chapitre~\ref{ch:transmission}). Du côté des bases elles-mêmes, A. E. Felle a présenté les problèmes et les perspectives de l’Epigraphic Database Bari, la principale base consacrée aux inscriptions des chrétiens de Rome\footnote{A. E. Felle, « Inscriptions by Christians in Late Antique Rome: Some Issues and Perspectives for the Epigraphic Database Bari (EDB) », dans I. Velázquez Soriano, D. Espinosa Espinosa (éd.), \emph{Epigraphy in the Digital Age}, Oxford, Archaeopress, 2021, p. 213–224 (notice vérifiée : catalogue de la TU München ; contenu non consulté).}.

\section{La chronologie consulaire}

Les \emph{Consuls of the Later Roman Empire} de R. S. Bagnall, A. Cameron, S. R. Schwartz et K. A. Worp (1987) sont l’instrument de référence pour l’usage des consulats comme système de datation. L’ouvrage documente, année par année, la reconnaissance différenciée des consuls en Orient, en Occident et dans les territoires des usurpateurs\footnote{R. S. Bagnall, A. Cameron, S. R. Schwartz, K. A. Worp, \emph{Consuls of the Later Roman Empire}, Atlanta, Scholars Press, 1987 (Philological Monographs of the American Philological Association, 36), viii–759 p., ISBN 1-55540-099-X (vérifié : Oxford University Press, FRANTIQ, HathiTrust). Comptes rendus : A. Chastagnol, \emph{Journal of Roman Archaeology}, 1, 1988, p. 203–206 ; J. R. Martindale, \emph{Journal of Roman Studies}, 79, 1989 ; J. Harries, M. Whitby, \emph{Classical Review}, 39, 1989.}. Les épitaphes de Rome y sont une source majeure, et des études récentes continuent de les exploiter pour l’histoire politique, par exemple pour l’usage romain des consuls de Magnence en 352\footnote{M. Humphries, « The Memory of Mursa: Usurpation, Civil War, and Contested Legitimacy Under the Sons of Constantine », dans N. Baker-Brian, S. Tougher (éd.), \emph{The Sons of Constantine, AD 337–361}, Cham, Palgrave Macmillan, 2020, p. 157–183, DOI 10.1007/978-3-030-39898-9\_6 (vérifié).}. La CLRE n’a pas été consultée directement. Elle est néanmoins citée à plusieurs reprises dans ce volume par l’intermédiaire des fiches de l’EDB, qui en reproduisent les pages pour les cas litigieux (p. 631, 636, 644, 661 et 680) : ces renvois sont toujours signalés comme indirects. La table des fastes du chapitre~\ref{ch:dater} reste à collationner sur la CLRE.

\section{La langue et les noms}

L’étude de la langue des inscriptions chrétiennes s’appuie sur les synthèses consacrées au latin tardif et « vulgaire », dont l’\emph{Introduction au latin vulgaire} de V. Väänänen et, pour la variation sociale, l’ouvrage de J. N. Adams (2013)\footnote{V. Väänänen, \emph{Introduction au latin vulgaire}, 3\textsuperscript{e}~éd., Paris, Klincksieck, 1981 (notice vérifiée) ; J. N. Adams, \emph{Social Variation and the Latin Language}, Cambridge, Cambridge University Press, 2013 (notice de l’éditeur vérifiée).}. L’onomastique chrétienne de Rome a été étudiée par I. Kajanto, dont les \emph{Onomastic Studies in the Early Christian Inscriptions of Rome and Carthage} (1963) et \emph{The Latin Cognomina} (1965) restent les références, et par Ch. Pietri (1977)\footnote{I. Kajanto, \emph{Onomastic Studies in the Early Christian Inscriptions of Rome and Carthage}, Helsinki, 1963 (Acta Instituti Romani Finlandiae, II, 1) ; \emph{The Latin Cognomina}, Helsinki, 1965 (notices vérifiées).}. Le chapitre~\ref{ch:noms} confronte le corpus à ces cadres, dans les limites d’un relevé automatique contrôlé.

\section{Les formules religieuses}

Aucune étude spécifiquement consacrée aux mentions de Dieu, du Christ, de l’Esprit ou de la Trinité dans les épitaphes chrétiennes de Rome n’a été identifiée. Le sujet est traité, à l’échelle du formulaire, dans les synthèses de Pietri (1983) et de Carletti (2008). Le présent volume aborde la question par le texte, en distinguant ce qui est gravé en lettres, abrégé, restitué, figuré par un signe ou seulement interprété (partie~VII). Cette approche ne remplace pas la confrontation avec ces synthèses, qui reste à faire sur pièces.
